% Financial Mathematics
% Exercise Template
%----------------------------------------------
\documentclass[addpoints,twoside]{exam}         % Can't specify font size here,
% so we'll set it manually below
\usepackage{problemset}                         % Custom package for problem sets -
% see problemset.sty
% -------------------------------------------------------------------------
\begin{document}

\FontEasyRead                                   % Set font size here

\Heading                                        % Modify heading as needed
{Financial Literacy}                            % Course
{3. Equivalent Cashflow Equations - Set 1}              % Title
{\textbf{Instructions: }For each problem, sketch a timeline, write an equation
in one variable, and solve for that variable to answer the question.}

\begin{questions}

  \Problem{An obligation of \$2,500 falls due at the end of 7 years.  If money is worth 10\% compounded
  monthly, find an equivalent debt at the end of 10 years.}

  \Problem{Hank owes \$1,000 due at the end of 18 months and \$1,500 due at the end of 4 years.  If
    money is worth 6\% compounded quarterly, what single payment in 2 years will liquidate
  both obligations?}

  \Problem{A debt of \$5,000 is due 5 years from now.  It is proposed instead that $\$X$ be paid now
    and another $\$X$ be paid 10 years from now.  Calculate the value of $X$ if the annual
  interest rate is 12\%.}

  \pagebreak

  \Problem[2.75in]{A man stipulates in his will that \$50,000 from his estate is to be placed in a fund
    from which his 3 daughters are each to receive \underline{equal} amounts when each reaches age 21.
    When the man dies, the girls are aged 19, 15, and 13.  If the fund earns 12\% compounded
  semiannually, how much will each daughter receive?}

  \Problem{John is repaying a debt with payments of \$250 per month at the end of each month.  If he misses his
    payments for July and August, what single larger payment in September will be required to put him back on
  schedule, if interest is 6\% compounted monthly?}

\end{questions}

\end{document}
